\section{Background and motivation}\label{sect:Introduction}

%%%%%%%%%%%%% Motivation  %%%%%%%%%%%%%%%%%%%%%

Since the work of Worthington \cite{worthington1881impact, worthington1883impact}, we have endeavoured to classify the different behaviours that arise when droplets impact a liquid surface. An important question is to distinguish under which conditions a droplet will rebound versus coalesce or splash \cite{rein1993phenomena}, especially in industrial processes such as spray injections \cite{moreira2010advances}.

The ability for a droplet to rebound depends entirely on whether the layer of air that separates the two fluids remains unbroken throughout the impact. In the event that this air layer becomes sufficiently thin, the Van der Waals forces between the two liquid regions act, and the liquid pierces through the air layer and initiates coalescence \cite{sprittles2024gas}. For rebound to occur, it is a delicate balance between the system having enough restoring force to jettison the drop upwards after a cushioned impact, but not so much that the air layer is breached due to the inertia of the drop before the liquid surface has time to force the droplet upwards. 
"inviscid fluid model for the water entry of a rigid, solid body into a quiescent pool

Industrial interest in solid impact dynamics were initially heavily motivated by military interest in sea-plane landings and slamming of boat hulls, leading to the development of Wagner theory \cite{wagner1932phenomena, korobkin2004analytical} as a foundational inviscid model for solid body entry into quiescent liquids. The role of the air layer in pre-impact dynamics of solid bodies on deep liquid baths was first developed by Smith et al. \cite{smith2003air}. For this work we base ourselves on the investigation of Moore \cite{moore2021introducing}, who combined pre-impact cushioning with the Wagner model. We incorporate surface tension effects in drop oscillations and in wave generation and propagation in the bath, relevant to scenarios such as spray coating and ink-jet printing (see Derby\cite{derby2010inkjet}) where surface tension dominates small droplet behaviour.. Additionally, studies of droplet-droplet impacts have been motivated by the rich field of aerosols in a wide variety of natural processes from cloud seeding \cite{bird2010daughter} to pathogen transfer through suspended droplets \cite{yang2007size}.


This study is originally motivated by phenomena where a vertical vibration of the liquid bath permits sustained repeated droplet bouncing to occur \cite{couder2005walking}. These Faraday pilot-wave systems, named due to the Faraday waves generated at sufficiently strong vibration, can present a wide range of complex behaviours. To obtain a strongly coupled drop-bath system, the vibration is held below that which destabilises the free surface (the Faraday threshold). Within this regime, there are multiple drop-bath behaviours that can be experimentally achieved and modelled, and have been quantified in work such as those of Protiere et al. \cite{protiere2006particle}, and Milewski et al. \cite{milewski2015faraday}. As the forcing reaches near the Faraday threshold, the bouncing modes are destabilised, and through the asymmetry of impact a horizontal force is applied to the drop, permitting them to `walk' along the surface of the bath, guided by the waves generated from previous bounces. This Faraday pilot-wave system has drawn parallels to the quantum theory described by de Broglie \cite{de1926interference}. Dubbed the hydrodynamic quantum analogue, attempts have been made to recreate macro-scale analogies to quantum phenomena. For comprehensive reviews see the works of Bush \cite{bush2015pilot, bush2018introduction} and references therein.
%%%%%%%%%%%%%%% Two Dimensional Behaviour %%%%%%%%%%%%%5
The nature of these studies has required long-time simulations of repeated bouncing, which quickly become computationally expensive. Efforts to minimise this cost include reducing the number of spatial dimensions, such as laid out in the work by Nachbin et al. \cite{nachbin2017tunneling}, where a two-dimensional system was used to investigate the relationship between a droplet and the bath topography and the resultant tunnelling effects between bath cavities, or linking coupled droplet dynamics \cite{nachbin2018walking}, and recent works have highlighted the importance of including vertical dynamics within these methods \cite{papatryfonos2024static}. 

Two-dimensional configurations provide comparatively fast, detailed understanding of the key behaviours which are computationally expensive in three-dimensions. This approach can also be viewed as the  `infinite-cylinder' approximation of the geometry as a two dimensional droplet represents the cross-section of a cylinder. This consideration in itself can be well-motivated by physical phenomena such as water-walkers, whose long feet are cylindrical to leading order \cite{hu2003hydrodynamics}.

%%%%%%%%%%%%% Droplet Dynamics %%%%%%%%%%%%%%%
The linear inviscid model for the oscillations of a droplet was developed in 1924 by Lamb\cite{lamb1924hydrodynamics}. The model has been validated experimentally by work such as in Becker et al. \cite{becker1991experimental}. Considering droplets in zero gravity cases allows for the uninterrupted study of their oscillatory nature \cite{wang1996oscillations} and the longer viscous damping timescales. Comparisons between computations of the Navier-Stokes equations and linear theory in two- and three-dimensions can be found in Aalilija et al. \cite{aalilija2020analytical}. Following Lamb, they used the inviscid solution and an energy argument to find linear damping rates. For bouncing droplets,  Mol\'{a}\u{c}ek \& Bush considered a simplified model where both bath and droplet deform equally as a logarithmic spring \cite{molavcek2013drops}. Terwagne et al. experimentally investigated the role of deformation in the droplet rebound dynamics when the drop and bath are in different regimes, and proposed a double mass-spring system to replicate this range of regimes in a drop-bath system \cite{terwagne2013role}. While solid spheres can rebound in this regime \cite{galeano2021capillary}, droplet deformation converts kinetic energy to drop deformations on contact with the bath, and then may revert it back to kinetic energy at take-off, increasing upwards propulsion.


%%%%%%%%%%%%%%%%% Regimes %%%%%%%%%%%%%%%%%%%%%%
Non-dimensional parameters play a key role in quantifying the relative importance of physical mechanisms in the flow. For rebounding droplets and the bath deformation, the three parameters of particular interest are the Weber, Reynolds, and Ohnesorge numbers. Which for a typical impact velocity $W_0$, length scale of the unperturbed drop radius $R_0$, dynamic viscosity $\mu$, surface tension $\sigma$, and density $\rho$, are given as

\begin{equation*}
   W\!e = \frac{\rho W_0^2 R_0}{\sigma}, \qquad Re = \frac{\rho W_0 R_0}{\mu} ,\qquad Oh = \frac{\mu}{\sqrt{\rho \sigma R_0}} =  \frac{\sqrt{W\!e}}{Re}, 
\end{equation*}

where here the Weber number indicates ratio of inertia to surface tension, the Reynolds number encapsulates the relative effect of inertia to viscosity, and the Ohnesorge number provides a relationship between viscosity and surface tension. There are several other parameters which could play a role; in particular for larger drops and longer waves the Bond number $Bo = \rho g R_0^2/\sigma$ gains importance. Even small changes to these parameters will affect the overall droplet behaviour, including whether the drops will bounce, or coalesce \cite{lewin2024collision}. The work presented within this paper will focus around the low-Weber and the quasi-inviscid low-Ohnesorge limit, in order to capture smooth rebound dynamics with small disturbances \cite{sprittles2024gas}.


%%%%%%%%%%%%%%%%% Current Literature %%%%%%%%%%%%
The investigation of single rebounds in this regime has thrived over recent years, with the main goal being computationally efficient models that faithfully reproduce the experimental and DNS results. 
The `kinematic match' method for axisymmetric \cite{galeano2017non} and non-axisymmetric impacts \cite{galeano2019quasi} modelled the droplet as a sphere, with the pressure on its contact surface obtained as a constraint for matching the shape of the sphere's and the surface of the bath. This method captured details of bouncing and walking droplets, matching experiments, and was later used in comparisons of solid sphere impacts \cite{galeano2021capillary}. The model neglected droplet deformation and the air layer. Alventosa \textit{et al.} \cite{alventosa2023inertio} use a simplified kinematic match methodology (matching the impact position at a single point) coupled with a deformable drop to model axisymmetric rebounds. They obtained good agreement with experiments and DNS, enabling a unified understanding of the bouncing dynamics in their target inertio-capillary regime (vanishing influence of viscosity and gravitational forces described by sufficiently small Ohnesorge and Bond numbers, i.e. $\text{Oh} \ll 1$ and $\text{Bo}\ll 1$).



%%%%%%%%%%%%%%%% SO WHAT %%%%
These methods do not attempt to capture the behaviour of the lubricating air layer during the rebound dynamics. In this work we seek a coupled system allowing for the inclusion of the lubricating air layer in the rebound dynamics of a two-dimensional drop off a deep liquid bath in the low-Weber, low-Ohnesorge regime. We derive a closed set of evolution equations that govern the evolutionary dynamics of the droplet, free surface together with a free-boundary elliptic problem for the air layer pressure. This methodology aims to fill the gap between the kinematic match methods and  DNS of the full fluid equations. The system presented herein can be simulated much more efficiently than full DNS, whilst still providing key insight into the influence of the air gap, as we will showcase in the sections to follow.

We next present the derivation of the system, considering each fluid region independently in Section \ref{sect:Model}. We then couple the regions together and present a numerical approach for simulating the results in Section \ref{sect:Numerics}. Finally, we will use DNS techniques that have been previously verified against experiments to determine the accuracy of our model, and demonstrate that it can produce efficient simulations of complex multiple-rebound dynamics off a vibrating bath in Section \ref{sect:Results} and Section \ref{sect:Faraday}, respectively. We conclude with a discussion of the model and consideration for future avenues of study in Section \ref{sect:Conclusion}.
